% === Section 6: Theorem Bridge ===
\section{Theorem Bridge: From Contracts to Safety}\label{sec:theory}

This section presents the formal theorem programme that connects the \dcs{} runtime
contracts to end-to-end safety guarantees.  The paper-facing presentation is
deliberately compact: four flagship theorem statements carry the argument, while
the permanent T1--T11 identifiers remain in place for registry, appendix, and
code-traceability purposes.  Definitions, corollaries, and propositions retain
their audit identifiers, but they are not presented as additional headline
theorems.  This avoids presenting the work as an inflated list of disconnected
theorems while preserving the full audit surface used by the repository.

\subsection{Headline Results and Audit Identifiers}\label{sec:headline-theorem-arc}

\begin{table}[h]
\caption{Main-paper theorem arc and permanent traceability identifiers.}
\label{tab:headline_theorem_arc}
\centering\small
\begin{tabularx}{\columnwidth}{@{}p{0.25\columnwidth}p{0.21\columnwidth}X@{}}
\toprule
\textbf{Headline result} & \textbf{Audit IDs} & \textbf{Role in the argument}\\
\midrule
Observation-aware safety gap & T1 & Degraded observation can make observed-state safety disagree with true-state safety.\\
Observation necessity & T4 & Observation-only or quality-ignorant mandatory release cannot guarantee safety on the defended ambiguity witness class.\\
ORIUS covered release guarantee & T2, T3 & Covered uncertainty sets plus repaired release yield one-step safety and an explicit risk-envelope corollary.\\
Typed transfer and covered optimality & T11; T10/T11 corollary & The typed ORIUS contract transfers to defended domains, with certificate/fallback discipline (T5--T8) and the covered Bayes lower/upper safety sandwich as scoped flagship or flagship-derived surfaces.\\
\bottomrule
\end{tabularx}
\end{table}

T5--T10 remain visible as scoped theorem surfaces rather than equal-size
headline claims.  The generated theorem-promotion matrix and result cards are
the authority for promoted status, proof anchors, code anchors, tests,
artifacts, and claim-boundary language.

\subsection{Assumptions}\label{sec:assumptions}

The theorem programme uses the canonical Appendix~B assumption register
$A1$--$A13$.  This condensed IEEE bridge restates the subset $A1$--$A9$
needed to present the battery ladder T1--T8.

\begin{assumption}[A1: Bounded Model Error]\label{ass:A1}
The one-step prediction error of the domain model is bounded:
$|\epsilon_t| \le \epsilon_{\mathrm{model}}$ for all $t$.
\end{assumption}

\begin{assumption}[A2: Bounded Telemetry Error]\label{ass:A2}
The observation error is bounded by a function of the reliability weight:
$|\hat{x}_t - x_t| \le g(w_t)$, where $g$ is a known, monotonically
non-increasing function with $g(1) = \bar\epsilon$ (nominal noise floor).
\end{assumption}

\begin{assumption}[A3: Constraint Expressibility]\label{ass:A3}
The domain safety specification can be expressed as a finite conjunction of
state-action constraints evaluable over $\Ut$.
\end{assumption}

\begin{assumption}[A4: Known Dynamics]\label{ass:A4}
The controller can evaluate the one-step dynamics map $f(x_t, a_t)$ to compute
$\Delta(a_t)$.
\end{assumption}

\begin{assumption}[A5: Monotone Bounded Inflation]\label{ass:A5}
The reliability-adaptive conformal margin $m_t$ is monotonically non-increasing in
$w_t$: lower reliability never shrinks the uncertainty set.  The base conformal
predictor achieves coverage $\ge 1 - \alpha$ on the calibration set.
\end{assumption}

\begin{assumption}[A6: Bounded Detector Lag]\label{ass:A6}
The staleness counter and drift detector respond within a bounded lag of
$\tau_{\mathrm{lag}} < \infty$ steps.
\end{assumption}

\begin{assumption}[A7: Causal Certificates]\label{ass:A7}
Each certificate uses only causally available information (no future data).
The episode horizon is finite: $T < \infty$.
\end{assumption}

\begin{assumption}[A8: Piecewise Battery Fallback Admissibility]\label{ass:A8}
In the battery witness row, fallback admissibility is piecewise: zero dispatch is
certified on the safe-set interior, a boundary-aware safe-landing recovery action
is certified near the boundary when feasible, and the runtime fails closed
otherwise.
\end{assumption}

\begin{assumption}[A9: Sub-Gaussian Drift Proxy]\label{ass:A9}
The battery certificate drift process admits a sub-Gaussian scale
$\sigma_{\mathrm{d}}$ and a boundary gap $\delta_{\mathrm{bnd},t}$, so the
certificate survival probability satisfies the closed-form first-passage lower
bound used by T6.
\end{assumption}

\begin{assumption}[A10: Boundary-Active Controller]\label{ass:A10}
The quality-ignorant baseline controller~$\pi$ used in the T1 witness construction
drives the true state within distance $\delta_{\mathrm{bnd}} = P_{\max}\Delta t$
of the safe-set boundary~$\partial\mathcal{S}$ with positive frequency.  This
captures the operational fact that economically-driven dispatch policies (e.g.\
deterministic arbitrage LPs in the battery domain and longitudinal
performance-tracking policies in the AV domain) operate near constraint
boundaries by construction; it is discharged empirically by the LP arbitrage
baseline used in~\S\ref{sec:battery} and the no-wrapper baseline in
\S\ref{sec:av}, both of which produce non-zero baseline TSVR under degraded
telemetry.  We state it as an assumption rather than a lemma because the
boundary-seeking property depends on the choice of controller, not on the
ORIUS framework itself.
\end{assumption}

\paragraph{Usage map.}  \Cref{tab:assumption_usage} shows which theorems depend on
which assumptions.

\begin{table*}[t]
\caption{Assumption usage map across the theorem programme.}
\label{tab:assumption_usage}
\centering\scriptsize
\resizebox{0.98\textwidth}{!}{%
\begin{tabular}{lcccccccccc}
\toprule
\textbf{Result} & \textbf{A1} & \textbf{A2} & \textbf{A3} & \textbf{A4} & \textbf{A5} & \textbf{A6} & \textbf{A7} & \textbf{A8} & \textbf{A9} & \textbf{A10} \\
\midrule
T1 OASG Existence         &\checkmark &\checkmark &  &\checkmark &  &  &  &  &  &\checkmark \\
T2 Safety Preservation    &\checkmark &\checkmark &\checkmark &  &\checkmark &  &\checkmark &  &  &  \\
T3 ORIUS Core Bound       &\checkmark &\checkmark &  &\checkmark &\checkmark &\checkmark &\checkmark &  &  &  \\
T4 Observation Necessity  &  &\checkmark &  &\checkmark &  &  &  &  &  &  \\
T5 Certificate Validity   &  &  &  &  &\checkmark &  &\checkmark &  &  &  \\
T6 Certificate Expiration &  &  &  &\checkmark &  &  &\checkmark &  &\checkmark &  \\
T7\_Battery Fallback      &\checkmark &  &  &\checkmark &  &  &  &\checkmark &  &  \\
T8 Graceful Degradation   &\checkmark &\checkmark &\checkmark &\checkmark &\checkmark &  &  &\checkmark &  &  \\
T9 Mandatory Release Impossibility&  &\checkmark &  &\checkmark &  &  &  &  &  &  \\
T10 Lower Bound           &  &\checkmark &  &  &\checkmark &  &  &  &  &  \\
T11 Typed Transfer        &  &  &\checkmark &\checkmark &\checkmark &  &\checkmark &\checkmark &  &  \\
\bottomrule
\end{tabular}}
\end{table*}

%% ----- Battery Ladder: T1--T4 -----
\subsection{Battery Ladder (T1--T4): The Reference Witness}

The battery domain serves as the \emph{reference witness}: it provides the deepest
closure of the proof chain because battery dynamics are scalar, linear, and fully
instrumented.

\begin{theorem}[T1: OASG Existence]\label{thm:T1}
Let $(\mathcal{X}, \mathcal{U}, f, \mathcal{S}, \hat{x})$ be a CPS satisfying
A1, A2, A4, and~\ref{ass:A10}.  Let $\pi$ be any quality-ignorant controller
satisfying~\ref{ass:A10} (boundary-active).  If the telemetry error satisfies
$\Prob[\|\xi_t\| > \delta_{\mathrm{bnd}}] > 0$, then there exists a time $t^*$
such that
\begin{equation}\label{eq:oasg_existence}
  \phi_{\mathrm{obs}}(\pi(\hat{x}_{t^*})) = 1
  \quad\text{and}\quad
  \phi_{\mathrm{true}}(\pi(\hat{x}_{t^*})) = 0.
\end{equation}
\end{theorem}

\begin{proof}[Proof sketch]
\emph{Step~1 (Boundary proximity).}
By~\ref{ass:A10}, the controller drives the true state within
$\delta_{\mathrm{bnd}} = P_{\max}\Delta t$ of the constraint boundary with
positive frequency.  Boundary-seeking is therefore a property of the controller
class, not a derived lemma; the assumption is discharged empirically by the
non-zero baseline TSVR observed for the LP arbitrage and no-wrapper policies in
\S\ref{sec:battery} and~\S\ref{sec:av}.

\emph{Step~2 (Observation gap).}
Under dropout rate $d > 0$, A2 (Bounded Telemetry Error) gives a non-degenerate
error distribution whose tail satisfies
$\Prob[\|\xi_t\| > \delta_{\mathrm{bnd}}] > 0$ for any positive
$\delta_{\mathrm{bnd}}$.

\emph{Step~3 (Witness construction).}
The boundary-proximity event of Step~1 and the tail event of Step~2 are
non-trivially overlapping under A4 (Known Dynamics) because the boundary event
has positive frequency and the tail event has positive probability per step.
At any step $t^*$ in the intersection, the observed safety predicate
$\phi_{\mathrm{obs}}$ evaluates to safe (the observed state is interior) but
the true safety predicate $\phi_{\mathrm{true}}$ evaluates to unsafe (the true
state has crossed the boundary).
\end{proof}

\paragraph{Battery instantiation.}
The true state is $x_t = \SOC_t \in [0, \Emax]$\,(MWh), the action is
$a_t = (P_t^{\mathrm{chg}}, P_t^{\mathrm{dis}})$, and the one-step dynamics are
\begin{equation}
  x_{t+1} = x_t + (\eta_c P_t^{\mathrm{chg}} - \eta_d^{-1} P_t^{\mathrm{dis}})\Delta t + \epsilon_t.
\end{equation}
The observed safety predicate evaluates the SoC constraint on $\hat{x}_t$:
$\phi_{\mathrm{obs}}(a_t) = \mathbf{1}[\SOC_{\min} \le \hat{x}_t + \Delta(a_t) \le \SOC_{\max}]$.

\paragraph{Empirical confirmation.}
Under the deterministic LP baseline with \texttt{drift\_combo} (20\% dropout),
the empirical TSVR is 3.9\% (DE) and 4.2\% (US CAISO), confirming
$\E[\TSVR] > 0$ for quality-ignorant controllers.

\begin{theorem}[T2: One-Step Safety Preservation Under the Shield]\label{thm:T2}
Under A1--A3, A5, and~A7, define the tightened safe-action set
\begin{equation}\label{eq:tightened_set}
  \Aset = \Bigl\{ a \in \mathcal{U} :
    f(\hat{x}_t, a) \oplus B(m_t) \subseteq \mathcal{S} \Bigr\},
\end{equation}
where $m_t = q_t / (w_t + \epsilon)$ is the reliability-inflated conformal margin
and $B(m_t)$ the ball of radius~$m_t$.  If $x_t \in \Ut$ and
$a_t \in \Aset$, then $f(x_t, a_t) \in \mathcal{S}$.
\end{theorem}

\begin{proof}[Proof (battery scalar case)]
Since $x_t \in \Ut = [\hat{x}_t - m_t, \hat{x}_t + m_t]$, we have
$|x_t - \hat{x}_t| \le m_t$.  The one-step dynamics give
\begin{align}
  s_{t+1} &= s_t + \eta_c P_t^{\mathrm{chg}}\Delta t - \eta_d^{-1}P_t^{\mathrm{dis}}\Delta t + \epsilon_t \notag\\
  &\ge \hat{s}_t - m_t + \eta_c P_t^{\mathrm{chg}}\Delta t
      - \eta_d^{-1}P_t^{\mathrm{dis}}\Delta t
      - \epsilon_{\mathrm{model}} \notag\\
  &\ge \SOC_{\min} + m_t - m_t - \epsilon_{\mathrm{model}} = \SOC_{\min} - \epsilon_{\mathrm{model}}.
\end{align}
Analogously for the upper bound.  Under A1, $\epsilon_{\mathrm{model}}$ is absorbed
into the margin, yielding $s_{t+1} \in [\SOC_{\min}, \SOC_{\max}]$.
\end{proof}

\begin{corollary}[T3: ORIUS Core Risk Envelope]\label{thm:T3}
Under A1, A2, A4--A7, the expected total violation count over horizon $T$ is bounded:
\begin{equation}\label{eq:core_bound}
  \boxed{\;\E[V_T] \;\le\; \alpha\,(1 - \bar{w})\,T\;}
\end{equation}
where $\alpha$ is the conformal miscoverage level and
$\bar{w} = T^{-1}\sum_{t=1}^{T} w_t$ is the episode-average reliability.
\end{corollary}

\begin{proof}
Define the per-step violation indicator $Z_t = \mathbf{1}[f(x_t, a_t^{\mathrm{rel}}) \notin \mathcal{S}]$.
By the calibration contract (A5), $\Prob[x_t \notin \Ut \mid \mathcal{H}_t] \le \alpha$.
Given the risk-envelope contract, the per-step risk budget is
$r_t = \alpha(1 - w_t)$, whence
$\Prob[Z_t = 1 \mid \mathcal{H}_t] \le \alpha(1 - w_t)$.  Summing:
\begin{align*}
  \E[V_T] = \sum_{t=1}^{T}\E[Z_t]
  &\le \sum_{t=1}^{T}\alpha(1-w_t) \\
  &= \alpha\sum_{t=1}^{T}(1-w_t)
   = \alpha(1-\bar{w})T. \qedhere
\end{align*}
\end{proof}

\paragraph{Numerical example.}
The numbers below are illustrative for the hypothetical moderate-degradation
row of \Cref{tab:core_bound_scenarios}, not the observed battery scenario.
With $\alpha = 0.10$, $T = 288$ (battery horizon), and a hypothetical
$\bar{w} = 0.83$ (moderate degradation):
$\E[V_T] \le 0.10 \times 0.17 \times 288 = 4.9$ violations.  The actual
battery runtime achieves $\bar{w} \in [0.905, 0.999]$ across all four
controller configurations (see \S\ref{sec:battery}), which corresponds to the
mild-to-perfect rows of the same table and predicts $\E[V_T] \le 2.74$
violations under the worst observed $\bar{w}=0.905$.  The empirically observed
violation count is 0, well below either bound.

\Cref{tab:core_bound_scenarios} illustrates the bound under varying degradation.

\begin{table}[h]
\caption{ORIUS core bound $\E[V_T] \le \alpha(1-\bar{w})T$ under varying
degradation scenarios ($\alpha = 0.10$, $T = 288$).}
\label{tab:core_bound_scenarios}
\centering\scriptsize
\begin{tabular}{lrrl}
\toprule
\textbf{Scenario} & $\bar{w}$ & $\E[V_T]$ & \textbf{TSVR (\%)} \\
\midrule
Perfect telemetry      & 1.00 & 0.0  & 0.00 \\
Mild degradation       & 0.95 & 1.4  & 0.50 \\
Moderate degradation   & 0.83 & 4.9  & 1.70 \\
Severe degradation     & 0.60 & 11.5 & 4.00 \\
Near-blackout          & 0.20 & 23.0 & 8.00 \\
Total blackout         & 0.00 & 28.8 & 10.00 \\
\bottomrule
\end{tabular}
\end{table}

\begin{theorem}[T4: Observation Necessity / No Free Safety]\label{thm:T4}
Under A2 and A4, for any non-trivial fault rate $d > 0$, there exists no
controller that simultaneously achieves $\TSVR = 0$ and zero intervention rate
(IR~$= 0$) across all episodes.
\end{theorem}

\begin{proof}[Proof sketch (constructive witness)]
Construct a degraded-observation episode where the deterministic LP drives SoC to
$\SOC_{\max} - \epsilon$ while the telemetry reads $\SOC_{\max} - \epsilon - \xi$
with $\xi > 0$.  Any controller that does not intervene will commit a dispatch that
pushes the true SoC above $\SOC_{\max}$.  Hence, maintaining $\TSVR = 0$ requires
IR~$> 0$.
\end{proof}

\subsection{Certificate and Temporal Theorems (T5--T8)}

\begin{theorem}[T5: Finite-Horizon Certificate Validity]\label{thm:T5}
Let $X_t$ be the observation-consistent state set at time~$t$, let
$\tilde a_{t:t+h}$ be the certified release or fallback sequence, and let
$S$ be the safe-state set.  If
\[
  \mathcal R_{t+j}(X_t,\tilde a_{t:t+j}) \subseteq S
  \qquad j=0,\ldots,h,
\]
then the certificate issued at time~$t$ is valid for horizon~$h$ unless a
certificate-invalidating event occurs.
\end{theorem}

\begin{proposition}[T6: Certificate Expiration Bound]\label{thm:T6}
Under A4, A7, and~A9, for any confidence level $\delta \in (0,1)$ the battery
certificate survival horizon satisfies the closed-form lower bound
\[
  \tau_t \;\ge\;
  \left\lfloor
    \frac{\delta_{\mathrm{bnd},t}}
    {\sigma_{\mathrm{d}}\sqrt{2\log(1/\delta)}}
  \right\rfloor.
\]
Runtime enforcement enters fallback when this lower bound drops below one step.
\end{proposition}

\begin{proposition}[T7\_Battery: Feasible Fallback Existence]\label{thm:T7}
Under A1, A4, and~A8, a constructive battery fallback exists in piecewise form:
zero dispatch is admissible and safe on the interior, a boundary-aware
safe-landing recovery action is admissible near the lower or upper boundary when
the actuator limits permit it, and the runtime fails closed otherwise.
\end{proposition}

\begin{theorem}[T8: Graceful Degradation Dominance with Useful Work]\label{thm:T8}
Let $P_G$ be the ORIUS graceful-degradation policy and $P_U$ an uncontrolled
persistence policy under the same admissible fault trace.  If
$Z_t^G \le Z_t^U$ for all~$t$ and $W_G\ge\lambda W_U$ for some
$\lambda\in(0,1]$, then $P_G$ weakly dominates $P_U$ in safety while preserving
at least a $\lambda$-fraction of useful work.
\end{theorem}

%% ----- Universal Extension: T9--T11 -----
\subsection{Universal Extension (T9--T11)}

The universal theorem is T11: any compliant domain inherits the ORIUS
runtime-assurance contract through the typed adapter obligations. T9 and T10
are promoted as scoped lower-bound surfaces: T9 covers mandatory
observation-only release with empty common safe core, and T10 covers a two-state
boundary-indistinguishability problem. Their status is governed by the
promotion matrix and result-card artifacts, not prose alone.

\begin{proposition}[T9: Impossibility of Quality-Ignorant Mandatory Release]\label{thm:T9}
For an observation $o$, define $B(o)=\{x:h(x)=o\}$ and
$K(o)=\cap_{x\in B(o)}C(x)$.  If $K(o)=\varnothing$, then no
observation-only mandatory-release controller $\pi(o)$ can guarantee
true-state safety for all $x\in B(o)$.
\end{proposition}

\begin{proof}
A mandatory-release controller chooses a single action $a=\pi(o)$.  Since the
common safe core is empty, no single action lies in every $C(x)$ for
$x\in B(o)$, so there exists $x^\dagger\in B(o)$ with $a\notin C(x^\dagger)$;
if the latent state realizes as $x^\dagger$, the released action is unsafe.
The result is a structural impossibility for the mandatory-release policy
class; abstention, fallback, uncertainty expansion, and denial are outside
this class by definition and are precisely the moves ORIUS adds.  The empty
safe-core hypothesis is discharged empirically by the non-zero baseline TSVR
of the LP arbitrage and no-wrapper policies in \S\ref{sec:battery} and
\S\ref{sec:av}: those policies operate as observation-only mandatory-release
controllers on ambiguity classes where $K(o)=\varnothing$ on a positive
fraction of steps.
\end{proof}

\begin{proposition}[T10: Boundary-Indistinguishability Lower Bound]\label{thm:T10}
Let $x_0,x_1$ be latent boundary states with induced observation
distributions $P_0,P_1$.  If $d_{\mathrm{TV}}(P_0,P_1)\le\epsilon$ and
$C(x_0)\cap C(x_1)=\varnothing$, then every observation-only mandatory-release
controller satisfies
\[
  \sup_{i\in\{0,1\}}\Pr_{O\sim P_i}\!\left[\pi(O)\notin C(x_i)\right]
  \ge \frac{1-\epsilon}{2}.
\]
\end{proposition}

\begin{proof}
Because the safe-action sets are disjoint, a single released action $\pi(O)$
must fail to lie in $C(x_i)$ for at least one $i\in\{0,1\}$ on any observation
realization.  Le~Cam's two-point inequality bounds the minimum sum of error
probabilities for any test between $P_0$ and $P_1$ from below by
$1 - d_{\mathrm{TV}}(P_0,P_1) \ge 1 - \epsilon$; taking the supremum over the
two latent labels yields the half of that lower bound stated above.
\end{proof}

\paragraph{Scope: TV distance is not the reliability weight.}
The bound is stated in terms of an exogenous TV distance $\epsilon$ between
two latent-state-induced observation laws.  $\epsilon$ is \emph{not} the ORIUS
reliability weight $w_t$.  The connection to the rest of the framework is
indirect: A2 (Bounded Telemetry Error) supplies a function $g(w_t)$ that
controls the observation error envelope, and a TV bound of the form
$\epsilon \le \epsilon(g(w_t))$ can be derived only for specific noise models
(Gaussian, bounded-support).  T10 therefore applies only to ambiguity classes
admitting such a TV bound; we do not claim that $w_t$ itself behaves as a TV
distance.

\paragraph{Reliability--risk frontier.}
T10 is a two-state lower bound, not a global minimax frontier.  The scoped
minimax row \texttt{T\_minimax} uses the same TV argument only over finite
two-state ambiguity classes and does not claim a matching universal upper bound.

\begin{theorem}[T11: Typed Structural Transfer]\label{thm:T11}
Let $\mathcal{D}$ be a domain instantiation satisfying the adapter contract
(\cref{sec:contracts}).  Assume that on every step $t$:
\begin{enumerate}[nosep,label=(\roman*)]
  \item \textbf{Coverage}: $\Prob[x_t \in \Ut \mid \mathcal{H}_t] \ge 1 - \alpha$.
  \item \textbf{Soundness}: Every $a \in \Aset(\Ut)$ keeps $f(x,a) \in \mathcal{S}$
        for all $x \in \Ut$.
  \item \textbf{Repair membership}: The repair operator returns
        $a_t^{\mathrm{safe}} \in \Aset(\Ut)$.
  \item \textbf{Fallback admissibility}: The domain certifies a safe hold or
        boundary-recovery fallback witness, and fails closed when neither case is
        admissible.
\end{enumerate}
Then:
\begin{equation}\label{eq:T11}
  \Prob\bigl[x_{t+1} \in \mathcal{S} \mid \mathcal{H}_t\bigr] \;\ge\; 1 - \alpha.
\end{equation}
\end{theorem}

\begin{corollary}[Covered Observation-Ambiguity Optimality]
\label{cor:T10_T11_observation_ambiguity_sandwich}
For an observation $o$, let
$B(o)=\{x:h(x)=o\}$ be its ambiguity class and
$K(o)=\bigcap_{x\in B(o)}C(x)$ be its common safe core.  Any mandatory-release
observation-only controller satisfies
\[
  \TSVR(\pi)
  \;\ge\;
  \mathbb{E}_{O}\!\left[
    \min_{a\in\mathcal{A}}
    \Prob(a\notin C(X^\star)\mid O)
  \right].
\]
If ORIUS constructs $\mathcal{X}_t$ with
$\Prob(X_t^\star\notin\mathcal{X}_t)\le\alpha$ and releases only
$u_t\in\bigcap_{x\in\mathcal{X}_t}C(x)$, then
$\TSVR_{\mathrm{ORIUS}}\le\alpha$; under exact coverage, the one-step
true-state violation probability is zero.
\end{corollary}

\begin{proof}
For each observation value $o$, define
$L(a,o)=\Prob(a\notin C(X^\star)\mid O=o)$.  A mandatory-release
observation-only policy must select one action $\pi(o)$ for the whole
ambiguity class $B(o)$, so
$L(\pi(o),o)\ge \min_{a\in\mathcal{A}}L(a,o)$.  Taking expectation over
$O$ gives the displayed Bayes lower bound.  If $K(o)=\emptyset$ and the
conditional law over $B(o)$ assigns positive mass to states excluding every
candidate action, the conditional minimum is positive; if $K(o)\ne\emptyset$,
the bound correctly permits zero conditional risk.

For ORIUS, split the violation event on whether
$X_t^\star\in\mathcal{X}_t$.  On the covered event, the released action is in
$C(x)$ for every $x\in\mathcal{X}_t$, hence in $C(X_t^\star)$.  Therefore the
covered contribution is zero, and the remaining contribution is bounded by
$\Prob(X_t^\star\notin\mathcal{X}_t)\le\alpha$.
\end{proof}

\paragraph{Counterexample boundary.}
The corollary is intentionally not a global optimality theorem.  It would be
false if a nonempty common safe core existed and an observation-only controller
selected from that core, if the controller were allowed to abstain instead of
mandatory release, if $\mathcal{X}_t$ failed to cover the true state, or if the
domain adapter supplied an incorrect safe set or infeasible fallback.  These
cases are treated as explicit proof obligations rather than hidden assumptions.

\paragraph{Failure-mode converse.}
If \emph{any} of the four T11 obligations fails, there exists a domain
instantiation where the transfer argument breaks and violations accumulate.

\paragraph{Episode aggregation corollary.}
If the domain instantiation supplies per-step risk budget $r_t$ with
$\Prob[x_{t+1} \notin \mathcal{S} \mid \mathcal{H}_t] \le r_t$, then
$\E[V_T] \le \sum_{t=1}^{T} r_t$.  With the battery envelope
$r_t \le \alpha(1 - w_t)$, this recovers the core bound: $\E[V_T] \le \alpha(1 - \bar{w})T$.

\subsection{Proof Architecture Summary}

The proof programme has a layered structure.  The main text should be read as
four flagship theorem statements plus supporting definitions, propositions, and
corollaries.  The T-identifiers serve as traceability anchors:
\begin{itemize}[nosep]
  \item \textbf{Problem} (T1): degraded observation creates the observed-safe /
        actually-unsafe gap.
  \item \textbf{Necessity} (T4): quality-ignorant mandatory release cannot
        guarantee true-state safety on the defended ambiguity witness class.
  \item \textbf{Construction} (T2, with T3 as a corollary): covered uncertainty
        sets, repair, and the risk-envelope contract supply the ORIUS upper
        guarantee.
  \item \textbf{Operationalization} (scoped T5--T8): certificate horizon, expiration,
        fallback, and graceful degradation make release time-bounded and
        fail-closed.
  \item \textbf{Transfer and optimality} (T11 plus the T10/T11 ambiguity
        corollary): typed obligations transfer the argument across defended
        domains, while the Bayes lower/upper sandwich states the covered
        ambiguity optimality claim.
\end{itemize}

The generated theorem-promotion matrix confirms that T9 and T10 are promoted
only in their scoped forms: mandatory-release empty-safe-core impossibility and
two-state boundary indistinguishability.  It is also the authority for proof,
code, test, artifact, hash, and claim-boundary status.
